\subsection{\texttt{TDIVS}}
\noindent\textbf{Bundle encoding.} UOP entry 19, opcode \texttt{0x101E}. Bundle: \texttt{B.OP + B.ARG + B.IOR + B.IOTI}. \texttt{B.OP.Opcode}=\texttt{0x101E}. \texttt{B.IOTI}: selector=\texttt{111 last}, \texttt{SrcTile0}=\texttt{src}, \texttt{SrcTile1}=\texttt{absent}, \texttt{DstTile}=\texttt{dst}, \texttt{SizeCode}=profile tile size. \texttt{B.IOR} imports scalar operand(s): \texttt{scalar}.\par\smallskip
{\small
\paragraph{PTO ISA description.}
Elementwise division with a scalar (tile/scalar or scalar/tile).
\paragraph{Example.}
Tile elements [10, 20, 30] divided by scalar 2 give [5, 10, 15].
\par\smallskip
\paragraph{PTO ISA operation.}
For each element \texttt{(i, j)} in the valid region:

\noindent{}Tile/scalar:

\[
\mathrm{dst}_{i,j} = \frac{\mathrm{src}_{i,j}}{\mathrm{scalar}}
\]

\noindent{}Scalar/tile:

\[
\mathrm{dst}_{i,j} = \frac{\mathrm{scalar}}{\mathrm{src}_{i,j}}
\]
}\par\smallskip
\paragraph{Notes.}
The scalar is imported via \texttt{B.IOR} and broadcast inside the microcode body. Division by zero produces Inf. Integer division truncates toward zero.
\paragraph{Microcode site table.}
{\scriptsize
\begin{longtable}{@{}R{0.055\linewidth}L{0.23\linewidth}L{0.31\linewidth}L{0.22\linewidth}@{}}
\toprule
Seq & Micro instruction & WO[63:32] fields & WM[31:0] fields \\
\midrule
0 & \texttt{u\_vec.vlds} & \texttt{opc=0x17, x=0, fl=mem, seq=0, kind=b} & \texttt{entry=19, cnt=2, res=vec, var=0} \\
1 & \texttt{u\_vec.vbr} & \texttt{opc=0x0B, x=0, fl=alu, seq=1, kind=c} & \texttt{entry=19, cnt=0, res=vec, var=0} \\
2 & \texttt{u\_vec.vdiv} & \texttt{opc=0x12, x=0, fl=alu, seq=2, kind=b} & \texttt{entry=19, cnt=2, res=vec, var=0} \\
3 & \texttt{u\_vec.vsts} & \texttt{opc=0x2A, x=0, fl=mem, seq=3, kind=b} & \texttt{entry=19, cnt=2, res=vec, var=0} \\
4 & \texttt{u\_vec.vlds} & \texttt{opc=0x17, x=0, fl=mem, seq=4, kind=b} & \texttt{entry=19, cnt=2, res=vec, var=0} \\
5 & \texttt{u\_vec.vbr} & \texttt{opc=0x0B, x=0, fl=alu, seq=5, kind=c} & \texttt{entry=19, cnt=0, res=vec, var=0} \\
6 & \texttt{u\_vec.vdiv} & \texttt{opc=0x12, x=0, fl=alu, seq=6, kind=b} & \texttt{entry=19, cnt=2, res=vec, var=0} \\
7 & \texttt{u\_vec.vsts} & \texttt{opc=0x2A, x=0, fl=mem, seq=7, kind=b} & \texttt{entry=19, cnt=2, res=vec, var=0} \\
\bottomrule
\end{longtable}
}
\paragraph{Microcode body DSL.}
\begin{lstlisting}[style=ptocode]
def uop_tdivs_tile_scalar(src, scalar, dst):
    dtype = src.element_type
    valid_rows, valid_cols = src.valid_shape

    for row in range(0, valid_rows, 1):
        remained = valid_cols
        for col in range(0, valid_cols, pto.get_lanes(dtype)):
            mask, remained = pto.make_mask(dtype, remained)
            vec = pto.vlds(src[row, col:])
            scalar_vec = pto.vbr(scalar)
            result = pto.vdiv(vec, scalar_vec, mask)
            pto.vsts(result, dst[row, col:], mask)
    return

def uop_tdivs_scalar_tile(scalar, src, dst):
    dtype = src.element_type
    valid_rows, valid_cols = src.valid_shape

    for row in range(0, valid_rows, 1):
        remained = valid_cols
        for col in range(0, valid_cols, pto.get_lanes(dtype)):
            mask, remained = pto.make_mask(dtype, remained)
            vec = pto.vlds(src[row, col:])
            scalar_vec = pto.vbr(scalar)
            result = pto.vdiv(scalar_vec, vec, mask)
            pto.vsts(result, dst[row, col:], mask)
    return
\end{lstlisting}
\Needspace{0.34\textheight}
